\chapter{Konklusjonen}

\section{Konklusjonen: fra forskningsspørsmålene til behandling av problemstillingen}

Etter å ha fullført de to rundene med spørreundersøkelse og gjennomført de fem dybdeintervjuene, har det blitt mulig å presentere et helhetlig svar på den sentrale problemstillingen i denne avhandlingen, basert på et integrert empirisk grunnlag.

De fire forskningsspørsmålene gjenspeiler en sammenhengende og sammenhengende kausal kjede; svaret på det første spørsmålet (RQ1) gjenspeilte at registreringsprosessen i den første versjonen manglet forutsigbarhet hos mange brukere, og dette skyldtes hovedsakelig uklarhet i interaksjonselementene. Med innføringen av de målrettede designendringene i den andre versjonen endret registreringsløpet seg, der det ble en enkel og tydelig opplevelse i øynene til alle deltakerne i intervjuene.

Når det gjelder det andre spørsmålet (RQ2), forklarer det den underliggende mekanismen bak denne endringen, og bekrefter at grensesnittets effektivitet i å informere brukeren om systemets status er svært viktig for å unngå manglende samsvar i den mentale modellen; dette skjedde da de strukturelle endringene på «bekreftelse av innsending, tydelighet i ikonene, og logikken for tilbakestilling ble justert» bidro til å redusere dette avviket på en merkbar måte. Intervjuene utvidet også forskningens kunnskapsmessige omfang ved å bevise at begrepet «tillit» gjenspeiler en flerdimensjonal tilstand som krever absolutt åpenhet i hele kjeden for databehandling, og ikke bare selve innsendingen.

Basert på dette identifiserte det tredje forskningsspørsmålet (RQ3) typer feil og knyttet dem til dimensjonene ved datakvalitet i henhold til følgende rekkefølge:

\begin{itemize}
    \item Uklarhet i den visuelle statusen: produksjon av dupliserte data (brudd på kriteriet for unikhet).
    \item Generering av fiktive data: bygging av uriktige koblinger (brudd på kriteriet for korrekt kobling).
    \item Tap av data under navigering: dannelse av ufullstendige registreringer (brudd på kriteriet for fullstendighet) \cite{wang1996}.
\end{itemize}

Intervjuene som ble gjennomført, la til en type nye strukturelle feil som viste seg i «feil kontekstuell klassifisering»; der brukeren blir tvunget til å velge en unøyaktig kategori på grunn av fraværet av det riktige alternativet som samsvarer med virkeligheten som burde være der. Dette gjenspeiler et direkte brudd på kriteriet for «egnethet» \cite{weichbroth2024} og en uunngåelig trussel mot kvaliteten på dataene som er beregnet for trening av kunstig intelligens-modeller \cite{heck2025, geiger2021}.

Derfor viste svaret på det fjerde spørsmålet (RQ4) muligheten for å løse disse problemene på to integrerte nivåer:

\begin{enumerate}
    \item Brukergrensesnittnivået (UI): der de tre gjennomførte designtiltakene (tekstetiketter, eksplisitt bekreftelse av innsending og logikk for tilbakestilling) bidro til å øke den samlede vurderingen av applikasjonen med 8,5 prosentpoeng.

    \item Når det gjelder det strukturelle og organisatoriske nivået: Intervjuene med aktørene fra selskapene (Stena og NG) og den uavhengige sektoren gjenspeilte at den operative tilpasningen av applikasjonen og dens egnethet for praktisk bruk nødvendigvis krever integrering av fleksible kontekstuelle kategorier, å oppnå programmatisk integrasjon med systemer for virksomhetsressursplanlegging (ERP), i tillegg til å arbeide med å aktivere funksjonen for midlertidig lagring og arbeid uten internettforbindelse. Disse resultatene fikk enighet og bekreftelse fra alle deltakerne i intervjuene.
\end{enumerate}

Disse resultatene gir et større empirisk bevis på at forholdet mellom grensesnittdesign, systemstatus og datakonsistens empirisk kan testes og måles kontekstuelt basert på forskningsmodellen \cite{nielsen1994,norman2013,wang1996}. Den metodiske trianguleringen som er brukt, gir mer pålitelighet til denne konklusjonen gjennom integrering av bevis; ettersom den kvantitative siden beviser at forbedring av visningen av systemstatus reduserer feilkildene som truer datakonsistensen, mens den kvalitative siden på den andre siden viser at bærekraftig utvikling av profesjonelle digitale plattformer krever en nøyaktig balanse mellom detaljene i det visuelle grensesnittet og den strukturelle tilpasningen til feltbaserte arbeidspraksiser.

Basert på det som er presentert, er den kjente tilnærmingen «dårlige inndata fører til dårlige utdata»  \cite{heck2025, geiger2021} ikke begrenset bare til fasen der algoritmene mates med data, men begynner fra starten av brukerens første interaksjon med innsamlingsgrensesnittet; påliteligheten her avhenger av en integrert dobbelthet: i hvilken grad grensesnittet tilbyr riktige alternativer som samsvarer med virkeligheten på den ene siden, og dets avhengighet av nøyaktigheten i brukerens egne valg på den andre siden.

Dette gir oss den viktigste lærdommen i denne studien: Visuell brukervennlighet alene er ikke nok som garanti for kvaliteten på treningsdata for kunstig intelligens. Den må heller knyttes til kontekstuell tilpasning, transparent behandling av dataflyten og full integrasjon med de eksisterende profesjonelle systemene, slik at applikasjonen kan oppnå en reell og bærekraftig operativ verdi i avfallshåndteringssektoren.